\clearpage
\section{Impacto ambiental}

Al realizar el análisis del impacto ambiental, se considera que la fabricación del dispositivo tendrá un impacto reducido, debido a que se optó por adquirir y ensamblar productos y piezas que ya se encuentran disponibles en el mercado. De esta manera, se busca disminuir la cantidad de procesos de manufactura necesarios para su construcción.

Al mismo tiempo, se utilizan materiales y componentes que se puedan obtener en las cercanías del Valle de México, para evitar envíos largos y disminuir la huella de carbono que estos trayectos generarían en el ambiente.

Se prevé que el consumo de energía eléctrica que demande el dispositivo sea algo alto, debido a la potencia consumida por las resistencias del sistema de temperatura y por la potencia demandada por el motor del eje de aireación. A su vez, se busca disminuir el impacto en el consumo de energía, optimizando las condiciones ambientales del contenedor, siendo, de esta forma, un recubrimiento térmico aislante una opción amigable con el medio ambiente, ya que reduce en gran medida el consumo eléctrico utilizado para mantener una temperatura constante en el sistema.

De esta forma, se busca generar un ambiente constante dentro del contenedor para que el uso de los actuadores de temperatura y aireación se utilicen lo menos posible, reduciendo el impacto negativo del sistema.

Por otra parte, siendo este un sistema que busca reducir el impacto ambiental generado por las plantas de compostaje tradicionales, eliminando la necesidad de transportar los residuos a un lugar especializado para su procesamiento y mejorando el aprovechamiento del producto obtenido.

Esto, en conjunto, nos genera un impacto ambiental menor y ayuda a disminuir una problemática ambiental que nos afecta hoy en día: el transporte y mal aprovechamiento de los residuos orgánicos.

En un principio, este proyecto tenía como objetivo procesar los residuos orgánicos generados en la cafetería escolar. Sin embargo, la cantidad de residuos producidos requería un sistema de mayores dimensiones, lo que implicaba un costo que excedía el presupuesto disponible para el proyecto. A pesar de ello, si fuera posible desarrollar un sistema con la capacidad de procesar la totalidad de estos residuos y replicarlo en cada escuela del Instituto Politécnico Nacional, se podría generar un impacto positivo significativo en la reducción de la contaminación y en el aprovechamiento de los residuos orgánicos generados en la Ciudad de México.